%% =====================================================================
%% IEEE Conference style final report
%% HOW TO USE (Overleaf):
%%   1. Open https://www.overleaf.com/latex/templates/ieee-conference-template/grfzhhncsfqn
%%   2. "Open as Template" -> replace its main.tex with THIS file
%%   3. Upload references.bib and your figures from results/ into the project
%%   4. Recompile. Target length: 10-12 pages excluding references.
%% Fill every [TODO] and replace placeholder numbers with your real results.
%% =====================================================================
\documentclass[conference]{IEEEtran}
\IEEEoverridecommandlockouts
\usepackage{cite}
\usepackage{amsmath,amssymb,amsfonts}
\usepackage{algorithmic}
\usepackage{graphicx}
\usepackage{booktabs}
\usepackage{textcomp}
\usepackage{xcolor}
\usepackage{url}

\begin{document}

\title{A Comparative Study of Classical Filters and Deep Learning for Image Denoising}

\author{\IEEEauthorblockN{[Your Name]}
\IEEEauthorblockA{Department of [Department]\\
Yuan Ze University\\
Email: [you@example.com]}}

\maketitle

\begin{abstract}
% 150-250 words. Problem, what you did, key result, conclusion.
Image denoising aims to recover a clean image from a noisy observation and is a
fundamental low-level vision task. This work presents a systematic comparison
between five classical denoising techniques---Gaussian, median, and bilateral
filtering, Non-Local Means (NLM), and Block-Matching and 3D filtering (BM3D)---and
a deep convolutional neural network, DnCNN, trained with residual learning for
blind Gaussian denoising. All methods are evaluated under identical conditions on
grayscale (Set12, BSD68) and color (Kodak) images, across additive white Gaussian
noise of varying strength and salt-and-pepper noise, using PSNR, SSIM, and runtime.
% [TODO] replace with your numbers:
Experiments show that DnCNN achieves the highest PSNR/SSIM on Gaussian noise
(e.g., XX.XX dB at $\sigma=25$ on Set12), narrowly surpassing BM3D while running
an order of magnitude faster at inference, whereas the median filter remains the
strongest baseline for salt-and-pepper noise. We analyze the trade-offs between
detail preservation, generalization across noise levels, and computational cost,
and discuss when classical priors remain competitive with learned models.
\end{abstract}

\begin{IEEEkeywords}
image denoising, DnCNN, BM3D, non-local means, convolutional neural network, PSNR, SSIM
\end{IEEEkeywords}

\section{Introduction}
% ~2 pages. Motivate denoising; noise model y = x + n; sensor/low-light context;
% classical -> learned evolution; contributions as a bullet list.
Digital images are inevitably corrupted by noise during acquisition and
transmission. Formally, a noisy observation is modeled as $y = x + n$, where $x$
is the latent clean image and $n$ the noise. [TODO: expand motivation.]

\noindent\textbf{Contributions.}
\begin{itemize}
\item A unified, reproducible benchmark comparing five classical filters with DnCNN
under identical noise and evaluation protocols.
\item Evaluation across grayscale and color images, multiple Gaussian noise levels,
and salt-and-pepper noise.
\item An analysis of accuracy vs.\ runtime trade-offs and the limits of generalization.
\end{itemize}

\section{Related Work}
% ~2 pages. Group by: spatial filters (Gaussian/median/bilateral \cite{tomasi1998}),
% non-local methods (NLM \cite{buades2005}, BM3D \cite{dabov2007}),
% learned methods (DnCNN \cite{zhang2017}, and successors). Cite the BSD dataset \cite{martin2001}.
Classical denoising spans local linear filters, edge-preserving filters such as the
bilateral filter \cite{tomasi1998}, and non-local methods including NLM
\cite{buades2005} and BM3D \cite{dabov2007}. Learning-based methods such as DnCNN
\cite{zhang2017} reframe denoising as residual regression. [TODO: expand.]

\section{Proposed Approach}
\subsection{Noise models}
We synthesize additive white Gaussian noise $n \sim \mathcal{N}(0,\sigma^2)$ for
$\sigma \in \{15,25,50\}$ and salt-and-pepper noise at 5\% density.

\subsection{Classical baselines}
We evaluate Gaussian, median, and bilateral filtering, NLM, and BM3D with the
parameter settings reported in Section~\ref{sec:exp}.

\subsection{DnCNN with residual learning}
DnCNN predicts the noise residual $\mathcal{R}(y)\approx n$ and recovers
$\hat{x} = y - \mathcal{R}(y)$. The network stacks Conv--BN--ReLU blocks
(17 layers, 64 feature maps). We train for \emph{blind} denoising by sampling
$\sigma \sim \mathcal{U}(0,55)$ per patch, minimizing
\begin{equation}
\mathcal{L}(\theta) = \frac{1}{2N}\sum_{i=1}^{N}\lVert \mathcal{R}_\theta(y_i) - (y_i - x_i)\rVert_F^2 .
\end{equation}

\begin{algorithmic}[1]
\STATE Extract $40\times40$ patches from Train400 (flip/rotate augmentation)
\FOR{each epoch}
  \FOR{each minibatch $\{x_i\}$}
    \STATE sample $\sigma$, form $y_i = x_i + n_i$
    \STATE $\hat{n}_i \leftarrow \mathcal{R}_\theta(y_i)$;\; update $\theta$ on $\lVert \hat{n}_i-n_i\rVert^2$
  \ENDFOR
\ENDFOR
\end{algorithmic}

% [TODO] add an architecture diagram exported from your slides:
% \begin{figure}[t]\centering\includegraphics[width=\columnwidth]{architecture.png}
% \caption{DnCNN architecture.}\label{fig:arch}\end{figure}

\section{Experimental Analysis}\label{sec:exp}
\textbf{Setup.} Training on Train400; testing on Set12/BSD68 (gray) and Kodak
(color). Metrics: PSNR, SSIM, average inference time. Implemented in PyTorch on a
Colab T4 GPU.

\begin{table}[t]
\caption{Grayscale denoising on Set12, $\sigma=25$. [TODO: fill from results/gray\_set12\_sigma25.csv]}
\label{tab:gray}
\centering
\begin{tabular}{lccc}
\toprule
Method & PSNR (dB) & SSIM & Time (s) \\
\midrule
Noisy (input) & XX.XX & 0.XXX & 0.00 \\
Gaussian      & XX.XX & 0.XXX & 0.XX \\
Median        & XX.XX & 0.XXX & 0.XX \\
Bilateral     & XX.XX & 0.XXX & 0.XX \\
NLM           & XX.XX & 0.XXX & 0.XX \\
BM3D          & XX.XX & 0.XXX & 0.XX \\
DnCNN (ours)  & \textbf{XX.XX} & \textbf{0.XXX} & 0.XX \\
\bottomrule
\end{tabular}
\end{table}

% Repeat similar tables for sigma=15/50 (CSV files) and color Kodak.
% Insert figures exported by the notebook:
% \begin{figure}[t]\centering\includegraphics[width=\columnwidth]{visual_gray.png}
% \caption{Visual comparison on Set12 at $\sigma=25$ (PSNR/SSIM annotated).}\label{fig:vg}\end{figure}

\subsection{Discussion}
% Address: which method wins where; detail vs. smoothing; salt-and-pepper (median best,
% DnCNN trained on Gaussian degrades); accuracy-vs-speed; color vs. gray; comparison
% to published DnCNN numbers \cite{zhang2017}.
[TODO]

\section{Conclusion}
We benchmarked classical and learned denoisers under a common protocol. [TODO: summarize
findings and future work: color-specific DnCNN, real-noise datasets, transformer denoisers.]

\section*{Data Availability Statement}
All source code, trained weights, result tables, and figures are available at:
\url{https://github.com/[your-username]/[your-repo]}. Datasets are publicly available
(BSD/Train400, Set12, BSD68, Kodak).

\bibliographystyle{IEEEtran}
\bibliography{references}

\end{document}
